\documentclass[11pt]{article}
\usepackage[T1]{fontenc}
\usepackage{textcomp}
\usepackage[margin=0.75in]{geometry}
\usepackage{amsmath,amssymb,amsthm}
\usepackage{array,booktabs}
\usepackage{hyperref}
\setlength{\parindent}{0pt}
\newtheorem{theorem}{Theorem}
\theoremstyle{definition}
\newtheorem{definition}{Definition}
\title{Flow Proof Artifact: prop-07-equiangular-implies-one-angle-proportional}
\author{Generated by \texttt{flow doc proof}}
\date{Flow Proof Book}
\begin{document}
\maketitle
If two triangles are equiangular, they have one angle equal and the sides about equal angles proportional.\\[0.5em]
\textbf{Source.} euclid — Elements, Book VI, Proposition 7\\[0.5em]
\bigskip
\noindent{\large \textbf{Derived fact 1.} \textit{If two triangles are equiangular, they have one angle equal and the sides about equal angles proportional} \hfill the Euclidean plane \cdot Euclid Book VI \cdot Proposition 7: equiangular triangles have one equal angle and proportional sides about it}\par
\smallskip\noindent\textit{Coordinate: the Euclidean plane \cdot Euclid Book VI \cdot Proposition 7: equiangular triangles have one equal angle and proportional sides about it}\par
\smallskip\noindent\textit{Source: euclid — Elements, Book VI, Proposition 7}\par
\smallskip\noindent\textit{Built on: proposition 6: one equal angle and proportional including sides imply equiangular triangles, for Euclid Book VI on the Euclidean plane}\par
\medskip
\noindent\textbf{Goal.} If two triangles are equiangular, they have one angle equal and the sides about equal angles proportional\par
\noindent$\displaystyle \text{similar implies one angle and sides}$\par
\medskip
\noindent
\begin{tabular}{@{} >{\raggedright\arraybackslash}p{0.06\textwidth} >{\raggedright\arraybackslash}p{0.47\textwidth} >{\raggedleft\arraybackslash}p{0.41\textwidth} @{}}
\textbf{\#} & \textbf{Proof} & \textbf{Mathematics} \\
\midrule
\textbf{1.} & We prove that proposition 7: equiangular triangles have one equal angle and proportional sides about it for Euclid Book VI the Euclidean plane. & \\[0.45em]
\textbf{2.} & Let ABC = triangle\_1. & \\[0.45em]
\textbf{3.} & Let DEF = triangle\_2. & \\[0.45em]
\textbf{4.} & Let angle\_B = angle\_E. & \\[0.45em]
\textbf{5.} & From step 2, step 3, and step 4, we can deduce that ratio(AB, BC) equals ratio(DE, EF). & $\displaystyle ratio(AB, BC) = ratio(DE, EF)$ \\[0.45em]
\textbf{6.} & We invoke the derived fact governing Euclid Book VI the Euclidean plane: proposition 6: one equal angle and proportional including sides imply equiangular triangles, for Euclid Book VI on the Euclidean plane. & \\[0.45em]
\textbf{7.} & From step 6, this implies similar implies one the conjunction of angle and sides. Hence proven. & $\displaystyle \text{similar implies one angle and sides}$ \\[0.45em]
\bottomrule
\end{tabular}
\label{thm:Geometry-Euclid Book VI-proposition_7_equiangular_triangles_have_one_equal_angle_and_proportional_sides_about_it}
\par\medskip\hrule
\end{document}
